\section{Conjugación en presente}
\textbf{Terminaciones regulares}
\begin{itemize}[leftmargin=*]
  \item ik: raíz (sin terminación).
  \item je/jij/u/hij/zij/het: raíz + \textit{-t}.
  \item wij/jullie/zij: infinitivo (raíz + \textit{-en}).
\end{itemize}

\textbf{Ejemplo con \textit{werken}}
\begin{center}
\begin{tabular}{lll}
\toprule
Persona & Forma & Ejemplo \\
\midrule
ik & werk & Ik werk thuis. \\
je/jij/u/hij/zij/het & werkt & Jij werkt in een café. \\
wij/jullie/zij & werken & We werken samen. \\
\bottomrule
\end{tabular}
\end{center}

\textbf{Práctica}
\begin{itemize}[leftmargin=*]
  \item Conjuga 5 verbos (\textit{wonen, leren, maken, spelen, koken}) en di una frase por persona.
  \item Escribe la tabla sin mirar y verifica.
\end{itemize}
